\subsection{分幂}

设 \(x \in M\) 且 \(k \in N\) 。显然 \(x_1 \otimes x_2 \otimes \ldots \otimes x_k\) ，其中

\[
x _ {1} = x _ {2} = \dots = x _ {k} = x,
\]

是 \(\mathbf{TS}^k (\mathbf{M})\) 的一个元素

定义 2. --- 若 \(x \in \mathbf{M}\) ，元素 \(x \otimes x \otimes \ldots \otimes x\) （\(\mathbf{T}\mathbf{S}^k(\mathbf{M})\) 的）用 \(\gamma_k(x)\) 表示。

命题 3. --- (i) 若 \(x \in \mathbf{M}\) ，则 \(p\) 次幂（\(x\) 的、在 \(\mathbf{T}\mathbf{S}(\mathbf{M})\) 中计算的那个）等于 \(p! \gamma_p(x)\) 。

\begin{enumerate}
\def\labelenumi{(\roman{enumi})}
\setcounter{enumi}{1}
\tightlist
\item
  设 \(x_{1},\ldots ,x_{n}\in \mathbf{M};\) 则
\end{enumerate}

\[
\gamma_ {p} \left(x _ {1} + x _ {2} + \dots + x _ {n}\right) = \sum_ {p _ {1} + p _ {2} + \dots + p _ {n} = p} \gamma_ {p _ {1}} \left(x _ {1}\right) \gamma_ {p _ {2}} \left(x _ {2}\right) \dots \gamma_ {p _ {n}} \left(x _ {n}\right).
\]

\begin{enumerate}
\def\labelenumi{(\roman{enumi})}
\setcounter{enumi}{2}
\tightlist
\item
  设 \(x_1, \ldots, x_n \in \mathbf{M}\) ，设 \(p_1, \ldots, p_n\) 为整数 \(\geqslant 0\) ，且 \(p = p_1 + \cdots + p_n\) 。设 \(E\) 为所有映射 \(\varphi\) （从 \(\{1, \ldots, p\}\) 到 \(\{1, \ldots, n\}\) 、满足下述条件者）组成的集合
\end{enumerate}

\[
\operatorname{Card} \varphi^ {- 1} (1) = p _ {1}, \dots , \operatorname{Card} \varphi^ {- 1} (n) = p _ {n}.
\]

则

\[
\gamma_ {p _ {1}} (x _ {1}) \gamma_ {p _ {2}} (x _ {2}) \dots \gamma_ {p _ {n}} (x _ {n}) = \sum_ {\varphi \in E} x _ {\varphi (1)} \otimes x _ {\varphi (2)} \otimes \dots \otimes x _ {\varphi (p)}.
\]

\begin{enumerate}
\def\labelenumi{(\roman{enumi})}
\setcounter{enumi}{3}
\tightlist
\item
  设 \(x \in \mathbf{M}\) ，并令 \(q, r\) 为整数 \(\geqslant 0\) 。则
\end{enumerate}

\[
\gamma_ {q} (x) \gamma_ {r} (x) = \frac {(q + r) !}{q ! r !} \gamma_ {q + r} (x).
\]

\begin{enumerate}
\def\labelenumi{(\alph{enumi})}
\setcounter{enumi}{21}
\tightlist
\item
  设 \(x_{1}, \ldots, x_{n} \in \mathbf{M}\) ，且对于 \(\mathrm{H} \subset \{1, \ldots, n\}\) 的记号，设 \(x_{\mathrm{H}} = \sum_{i \in \mathrm{H}} x_{i}\) ，则
\end{enumerate}

\[
(- 1) ^ {n} x _ {1} x _ {2} \dots x _ {n} = \sum_ {\mathrm{H} \subset \{1, \dots , n \}} (- 1) ^ {\text { Card   H }} \gamma_ {n} (x _ {\mathrm{H}}).
\]

断言 (i) 直接由命题 2 (ii) 得出。

我们来证明 (ii)；对 \(n\) 作归纳可知，只需考虑 \(n = 2\) 的情形。于是我们有

\[
\begin{array}{l} \gamma_ {p} (x _ {1} + x _ {2}) = (x _ {1} + x _ {2}) \otimes (x _ {1} + x _ {2}) \otimes \dots \otimes (x _ {1} + x _ {2}) \quad (p \text {factors}) \\ = \sum_ {p _ {1} + p _ {2} = p} \sum_ {\sigma \in \mathfrak {S} _ {p _ {1}, p _ {2}}} \sigma (x _ {1} \otimes x _ {1} \otimes \dots \otimes x _ {1} \otimes x _ {2} \otimes x _ {2} \otimes \dots \otimes x _ {2}) \\ = \sum_ {p _ {1} + p _ {2} = p} \sum_ {\sigma \in \mathfrak {S} _ {p _ {1}, p _ {2}}} \sigma (\gamma_ {p _ {1}} (x _ {1}) \otimes \gamma_ {p _ {2}} (x _ {2})) \\ = \sum_ {p _ {1} + p _ {2} = p} \gamma_ {p _ {1}} (x _ {1}) \gamma_ {p _ {2}} (x _ {2}). \end{array}
\]

为了证明 (iii)，设 \(\mathfrak{S}_{p_{1},\ldots,p_{n}}\) 为 \((1,p_{1}+\cdots+p_{n})\) 中满足下列条件的置换组成的集合：它们在区间

\[
\left(1, p _ {1}\right), \left(p _ {1} + 1, p _ {1} + p _ {2}\right), \dots , \left(p _ {1} + \dots + p _ {n - 1} + 1, p _ {1} + \dots + p _ {n}\right)
\]

上的限制是递增的。根据 I，第 60 页，例 2 及命题 2 (ii)，我们有

\[
\gamma_ {p _ {1}} \left(x _ {1}\right) \gamma_ {p _ {2}} \left(x _ {2}\right) \dots \gamma_ {p _ {n}} \left(x _ {n}\right) = \sum_ {\rho \in \mathfrak {S} _ {p _ {1}, p _ {2}, \dots , p _ {n}}} \rho \left(x _ {1} \otimes x _ {1} \otimes \dots \otimes x _ {1} \otimes x _ {2} \otimes x _ {2} \otimes \dots \otimes x _ {2} \otimes \dots \otimes x _ {n} \otimes x _ {n} \otimes x _ {n} \otimes \dots \otimes x _ {n}\right)
\]

（其中有 \(p_i\) 个因子 \(x_i\) ）且此和等于

\[
\sum_ {\varphi \in E} x _ {\varphi (1)} \otimes x _ {\varphi (2)} \otimes \dots \otimes x _ {\varphi (p)}.
\]

在 (iii) 中令 \(n = 2\) ，\(x_{1} = x_{2} = x\) ，\(p_{1} = q\) 且 \(p_{2} = r\) ，则我们得到 (iv)（I，同上）。

最后，(v) 由命题 2 (ii) 以及 I 第 100 页命题 2 应用于元素 \(x_{i}\) （环 \(\mathbf{T}(\mathbf{M})\) 的）得出。

注记。--- 1) 设 \((x_{i})_{i\in I}\) 是 M 中元素的一个族。对于每个 \(\nu \in \mathbf{N}^{(I)}\) 的记号，设

\[
x _ {\nu} = \prod_ {i \in I} \gamma_ {\nu_ {i}} (x _ {i}).
\]

若 \((\lambda_{i})\in \mathbf{A}^{(\mathrm{I})}\) 且 \(p\in \mathbf{N}\) ，则由命题 3 (ii) 我们有

\[
\gamma_ {p} \left(\sum_ {i \in I} \lambda_ {i} x _ {i}\right) = \sum_ {\nu \in \mathbf {N} ^ {(I)}, | \nu | = p} \lambda^ {\nu} x _ {\nu}.\tag{6}
\]

\begin{enumerate}
\def\labelenumi{\arabic{enumi})}
\setcounter{enumi}{1}
\tightlist
\item
  设 \(\mathcal{M}\) 为从 \([1, p]\) 到 I 中的映射组成的集合。我们定义一个映射 \(\rho \mapsto \rho^{*}\) （从 \(\mathcal{M}\) 到 \(\mathbf{N}^{(I)}\) ），令
\end{enumerate}

\[
\rho^ {*} (i) = \operatorname{Card} \rho^ {- 1} (i).
\]

两个元素 \(\rho_1, \rho_2\) （\(\mathcal{M}\) 的）满足 \(\rho_1^* = \rho_2^*\) ，其必要且充分条件是存在 \(\sigma \in \mathfrak{S}_p\) 使得 \(\rho_2 = \rho_1 \circ \sigma\) （I，第 95 页）。由命题 3 (iii)，对于 \(|\nu| = p\) ，我们有

\[
x _ {\nu} = \sum_ {\rho \in \mathcal {M}, \rho^ {*} = \nu} x _ {\rho (1)} \otimes x _ {\rho (2)} \otimes \dots \otimes x _ {\rho (p)}.\tag{7}
\]

